\section{Quadratic systems and the certificate question}
\label{sec:setting}

Let $n,m\geq 1$, and let
\begin{equation}\label{eq:system}
 f_i(x)=x\transpose A_i x+2b_i\transpose x+c_i,
 \qquad
 S=\{x\in\R^n:f_i(x)<0\text{ for }i=1,\ldots,m\},
\end{equation}
where $A_i\in\Sym^n$, $b_i\in\R^n$, and $c_i\in\R$.
Here $\Sym^n$ denotes the real symmetric $n\times n$ matrices;
$A\succeq0$ and $A\succ0$ mean positive semidefinite and positive
definite, respectively. We use the Euclidean vector norm and the
Frobenius matrix norm, and write $\langle A,B\rangle=\tr(AB)$.
The symbol $\conv S$ denotes the ordinary convex hull, without taking
its closure. A hull is \emph{proper} if it differs from $\R^n$.
Our question is whether a nonempty quadratic system has a proper hull,
and how this can be certified by its defining inequalities.

For $\lambda\in\R^m$, write
\[
 (A_\lambda,b_\lambda,c_\lambda)
   =\sum_{i=1}^m\lambda_i(A_i,b_i,c_i),
 \qquad f_\lambda=\sum_{i=1}^m\lambda_i f_i.
\]
A nonzero $\lambda\geq0$ defines an \emph{aggregation}
$f_\lambda(x)<0$, which is valid for every $x\in S$.
We call $\lambda$ a \emph{convex certificate} if $A_\lambda\succeq0$,
so that $f_\lambda$ is convex on $\R^n$.
It is \emph{nontrivial} if $(A_\lambda,b_\lambda)\ne(0,0)$.
Otherwise the aggregated function is constant, and if $S\ne\varnothing$
evaluation at a point of $S$ shows that the constant is strictly
negative. For convenience the following cone contains zero:
\begin{equation}\label{eq:certificate-cone}
 \mathcal K=\{\lambda\in\R^m_+:A_\lambda\succeq0\},
 \qquad
 \mathcal T=\{\lambda\in\R^m:A_\lambda=0,\ b_\lambda=0\}.
\end{equation}
Thus a nontrivial convex certificate exists precisely when
$\mathcal K\not\subseteq\mathcal T$.

The homogeneous matrices and map associated with~\eqref{eq:system} are
\begin{equation}\label{eq:homogenization}
 Q_i=\begin{pmatrix}A_i&b_i\\b_i\transpose&c_i\end{pmatrix},
 \qquad
 f_i^h(x,t)=x\transpose A_i x+2t b_i\transpose x+c_i t^2,
 \qquad f^h=(f_1^h,\ldots,f_m^h).
\end{equation}
We also write $Q_\lambda=\sum_i\lambda_iQ_i$ and
$g_i(x)=x\transpose A_i x$.
Define the strict homogeneous feasible set
$S^h=\{(x,t):f_i^h(x,t)<0\text{ for all }i\}$.
It is open and invariant under multiplication by every nonzero real
scalar. For $t\ne0$, $(x,t)\in S^h$ is equivalent to $x/t\in S$.
For $\alpha\in\R^n\setminus\{0\}$ and $s\in\R$, put
\begin{equation}\label{eq:hyperplanes}
 H_{\alpha,s}=\{(x,t)\in\R^{n+1}:\alpha\transpose x=st\},
 \qquad E=\{(x,t):t=0\}.
\end{equation}

\begin{definition}[Hidden hyperplane convexity]
\label{def:hhc}
A homogeneous quadratic map $\phi:\R^d\to\R^m$ has
\emph{hidden hyperplane convexity} (HHC) if $\phi(H)$ is convex for
every linear hyperplane $H\subset\R^d$.
\end{definition}

This is the definition of Blekherman, Dey, and Sun
\citep[Definition~2.1]{BDSpreprint}, published as \citep{BDS2024}. Ordinary \emph{hidden convexity}
means only that $\phi(\R^d)$ is convex; the two properties differ.
Every homogeneous quadratic image contains zero and is invariant under
nonnegative scaling, so a convex such image is a convex cone, not
necessarily closed.
HHC implies ordinary hidden convexity when $d\geq3$:
any two points of $\R^d$ lie in some linear hyperplane, whose image
contains the segment joining their images.

\begin{question}[Blekherman--Dey--Sun, Conjecture~3.3]
\label{question:bds}
Suppose $S\ne\varnothing$ and $f^h$ has HHC. Is
\begin{equation}\label{eq:bds-question}
 \conv S\ne\R^n
 \quad\Longleftrightarrow\quad
 \exists\lambda\in\R^m_+:
 A_\lambda\succeq0,\quad (A_\lambda,b_\lambda)\ne(0,0)?
\end{equation}
\end{question}

The formulation in \citet[Conjecture~3.3]{BDSpreprint} says that the
aggregated function is not a negative constant; under nonemptiness this
is equivalent to the displayed condition. Their standing setting has
$n\geq3$ and $m\geq2$; ours includes all $n,m\geq1$. Numbered
locators for that work refer to the May~29, 2023 arXiv version; the
earlier author manuscript numbers several background results
differently.

The reverse implication in~\eqref{eq:bds-question} is elementary:
a nonconstant convex quadratic aggregation bounds $S$ by a proper convex
sublevel set. Under convexity of $(g_1,\ldots,g_m)(\R^n)$, Blekherman,
Dey, and Sun prove that $\mathcal K=\{0\}$ implies $\conv S=\R^n$
\citep[Proposition~2.14]{BDSpreprint}, and for $n\geq2$ their separable
case, with all $Q_i$ diagonal, gives the certificate characterization
\citep[Theorem~2.12]{BDSpreprint}. The unresolved case is the forward
implication when $\mathcal K$ contains nonzero vectors but all of them
belong to $\mathcal T$.

The certificate question differs from a complete description of the
hull. In the aggregation literature, a
\emph{good aggregation} is a nonnegative nonzero multiplier $\lambda$
such that $\conv S\subseteq\{x:f_\lambda(x)<0\}$ and $Q_\lambda$
has at most one negative eigenvalue. The function $f_\lambda$ need not
be globally convex. Blekherman, Dey, and Sun prove that these good
aggregations describe $\conv S$ under HHC, $n\geq3$,
and $\varnothing\ne S$ with proper hull
\citep[Theorem~2.9]{BDSpreprint}.

The two-constraint case is classical. Dines' theorem gives convexity of
the joint image of two homogeneous quadratic forms \citep{Dines1941},
and therefore HHC for any pair of homogeneous forms and their
restrictions. Yildiran describes the hull of a nonempty set defined by
two strict quadratic inequalities using at most two aggregations
\citep{Yildiran2009}. For three quadratic inequalities in dimension
$n\geq3$ with a nonempty proper hull, Dey, Mu\~noz, and Serrano obtain
an aggregation description when $Q_1,Q_2,Q_3$ admit a positive definite
real linear combination, with coefficients of either sign
\citep{DMS2022}. Blekherman and Dunbar use the spectral curve and
algebraic topology to obtain further three-quadratic hull descriptions
and finiteness results \citep{BD2025,BDpreprint}. Neither Dines'
theorem nor the elementary separation and compactness facts used below
are contributions of this paper.
